This thesis presents the development of an adaptive traffic signal controller based on \textbf{Deep Reinforcement Learning (DRL)} using the \textbf{Proximal Policy Optimization (PPO)} algorithm integrated with the microscopic traffic simulator \textbf{SUMO}. The main objective is to reduce urban congestion by dynamically optimizing signal timings in a high-demand intersection.

The agent learns to manage highly variable traffic conditions through a reward function designed to balance efficiency and stability, incorporating penalties for queue length, directional imbalance, and frequent phase switching. Training was conducted under a dynamic 60-minute traffic profile with several demand peaks.

Results show an \textbf{approximate 60\% reduction in average queue length} and a \textbf{4.5\% increase in throughput} compared to a fixed-time controller. These findings demonstrate the effectiveness of combining DRL and SUMO to improve urban mobility. Limitations and future work directions are discussed, including multi-agent control and perception via computer vision.

\vspace{1em}
\noindent \textbf{Keywords:} intelligent traffic lights, reinforcement learning, PPO, SUMO, urban mobility.